% Sections II-V for the ICSE 2027 tool paper. Drop-in after the Introduction.
% \toolname is the tool's name (the Introduction currently calls both the
% method and the tool "REAL"; a distinct tool name avoids that ambiguity).
% Every claim below is implemented; results the study has not produced yet are
% marked \todo{} and must not be filled with estimates.
\providecommand{\toolname}{REAL}
\providecommand{\todo}[1]{\textcolor{red}{[TODO: #1]}}

\section{Background}
\label{sec:background}

REAL~\cite{real2026} builds on Zave and Jackson's entailment
$D, S \models R$~\cite{zave1997}: a system specification $S$, under domain
assumptions $D$, must satisfy the requirement $R$. A failure is therefore
evidence against $S$ only if $D$ held when it occurred. REAL calls such a
failure \emph{valid}; a failure in a situation that breaks $D$, i.e.\ outside
the operational design domain (ODD), is \emph{spurious}. Valid failures are
read as \emph{obstacles} in the KAOS sense~\cite{lamsweerde2000}: conditions
under which a goal cannot be satisfied. REAL iterates
\emph{identify--analyse--mitigate}: a grammar generates test scenarios from
the requirement, failures are grouped into obstacles, and each obstacle is
resolved at one of four layers: the data, the model, the system around the
model, or the requirement itself. REAL describes this loop; it gives no
tool support for running it.

\section{\toolname}
\label{sec:tool}

\toolname{} supports one round of the loop in four stages
(Fig.~\ref{fig:arch}). The engineer writes the requirement and decides;
everything else runs automatically, and each computed result is shown with
the evidence behind it. The stages exchange four artefacts: the checked
requirement, the execution traces, the diagnosis, and the signed decisions.

\subsection{Stage 1: Requirement Elicitation}
\label{sec:specify}
Requirements are written in the REAL DSL~\cite{real2026}, which
\toolname{} extends with explicit slots for the three parts of the entailment:
\begin{itemize}
\item a goal refined into tasks, each \texttt{performed by} a named module
(the system $S$, e.g.\ a detector and a braking controller);
\item \texttt{assuming} clauses (the domain $D$), e.g.\
\texttt{"fog\_density <= 50"}, written over quantities the simulator can
measure, or as free text that is recorded but not checked;
\item \texttt{ensuring} clauses (soft goals in $R$), e.g.\
\texttt{"resume\_within\_s <= 10"}.
\end{itemize}
The \emph{requirement editor} presents the requirement as text with gaps.
Bounded choices are menus read from the code: the modules the executor
implements, the measurable quantities, the operators. The \emph{REAL-DSL parser}
checks every edit and marks the line it rejects, and the requirement can be
viewed as a goal tree or a KAOS diagram. The \emph{requirement check} then
applies two refusals that keep the entailment honest. An assumption about the
system under test (any quantity of the car itself) is rejected, since it
would let any failure be defined away. A misspelt or unimplemented module
name is rejected instead of silently running a different system. The stage
yields $R_0$, checked, with its goal model.

\subsection{Stage 2: Generate and Execute Tests}
\label{sec:generate}
A BNF grammar defines the scenario space from the requirement; its factors
are the pedestrian, clothing, direction of crossing, fog, time of day, road
wetness and an occluding parked car. The \emph{scenario search} enumerates
the space exhaustively (a grid) or explores it by grammatical
evolution~\cite{oneill2001,grape2022}, which deliberately goes beyond the
stated assumptions. Each scenario becomes a Scenic program~\cite{scenic2023}
executed in CARLA~\cite{carla2017} with the system under test in the loop:
the named detector and braking controller. A \emph{safety monitor} scores
every run and a \emph{telemetry recorder} writes its per-step trace:
distance to the pedestrian, car and pedestrian speed, detector confidence,
and the pedestrian's position relative to the lane. The stage runs as a job
on a GPU cluster. Before any cluster time is spent, two checks run: every
scene is compiled on the laptop against the real map, and the cluster's copy
of the code is compared with the local one. The run panel follows the job,
shows partial results as simulations finish, and retrieves the execution
traces when it ends.

\subsection{Stage 3: Diagnose Test Results}
\label{sec:diagnose}
The diagnosis reads the execution traces, never the simulator. The outcome
of each run is computed from its trace under a path-aware safety rule: the
car must keep a 5\,m margin while the pedestrian is in its lane. The same
function scores the run in the cluster's monitor, in the search's fitness
and here, so a search cannot exploit a rule the analysis does not use. Three
components turn the runs into a diagnosis.
\begin{itemize}
\item \emph{Validity check.} For every run, each \texttt{assuming} clause is
\emph{held}, \emph{broken} or \emph{not measured}, and a failure is spurious
only if a clause was broken; each \texttt{ensuring} clause is met or missed.
Runs in which the test itself did not happen (no encounter, the car left the
road, the pedestrian never entered the lane) are set aside and listed as test
defects, never counted as passes.
\item \emph{Failure classifier.} Each valid failure is typed by what its
trace shows: never detected, detected but not braked, detected too late,
brake released, braking insufficient, or stopped too close. Passes in which
the car stopped and never drove on are kept as their own type, since a car
that freezes is safe yet still misses its soft goal.
\item \emph{Obstacle analyser.} Obstacles from a catalogue (e.g.\
\textsc{AdverseWeather}, \textsc{DetectionTooLate}) and unnamed candidates
are retained only when a scenario factor raises the failure rate by at least
15 points with at least 10 failures. Each obstacle is attached to the goal it
blocks in the goal model and comes with its failure rates by factor level,
example traces, and candidate mitigations per layer. Each assumption also
receives a verdict: \emph{load-bearing} (failures rise when it is broken),
\emph{not load-bearing}, \emph{fewer failures when broken}, \emph{untested},
or \emph{insufficient data}, so the question \emph{``was this assumption
needed?''} is answered from the data.
\end{itemize}
A map of behaviours complements the diagnosis: every valid encounter is
placed by an edit distance over its sequence of events (multidimensional
scaling with hierarchical clustering); the map reports its own faithfulness
and silhouette, and labels weak clusters as neighbourhoods, not types.

\subsection{Stage 4: Decide and Refine}
\label{sec:decide}
In the \emph{review interface} the engineer decides each obstacle:
\emph{accept}, \emph{rename}, \emph{reject} or \emph{defer}, with a reason
in their own words; chooses, for each accepted obstacle, a mitigation from a
menu spanning the data, model, system and requirement layers; and rules on
each assumption: keep, tighten, loosen, drop, or add one. The tool's verdicts
are suggestions only; nothing is decided without the engineer. From the decisions, the interface writes $R_1$ next to $R_0$ with every
change labelled by the part of the entailment it touches: [S] the
system, [D] an assumption, [R] the requirement, or [T] the test itself (a
defect in the scene, which is fixed rather than blamed on the system). The
engineer signs, which records the signed decisions: each verdict, reason and
mitigation, the reviewer, and the run they rest on; a requirement that does
not parse or cannot be executed as written cannot be signed. The labels
decide what happens next (the diamond in Fig.~\ref{fig:arch}): if the system
changed ([S]), the next round starts from stage~1 with $R_1$ as its input; if
only assumptions or the requirement changed ([D], [R]), stage~3 re-diagnoses
the stored traces in seconds, with no new simulation. Each new run records
its parent run and the decision that caused it, so the lineage view shows how
a requirement evolved and why.

\subsection{Implementation}
\label{sec:impl}
The back end is written in Python: the DSL parser, a FastAPI service that
drives Scenic and CARLA, the analysis, and the job control (Slurm with an
Apptainer image). The front end is a React/TypeScript single-page
application. Every result is a file in a run folder (metadata with code and
grammar hashes, per-simulation tables, traces), so any analysis can be
recomputed from the traces without the simulator. A Docker image runs the
complete review loop on recorded rounds without a GPU or simulator. The tool
has \todo{N} automated tests, none of which needs the simulator.

\section{Preliminary Evaluation}
\label{sec:eval}

We ask: \textbf{RQ}: \emph{How effective is \toolname{} in helping engineers
validate and refine autonomous-system requirements and domain assumptions?}
Engineers may legitimately disagree on what a failure means, so we measure
what the diagnoses lead people to do, and where they diverge, rather than
agreement with a single answer.

\textbf{Subject.} An automated-braking system (a YOLOv5 detector and a
proportional braking controller) under a requirement assuming
\texttt{fog\_density <= 50}, \texttt{sun\_altitude\_angle >= 0} (no night),
an initial separation of at least 15\,m and a pedestrian speed of at most
3\,m/s, and ensuring that the car drives on within 10\,s. The grammar
enumerates 288 scenarios (pedestrian $\times$ clothing $\times$ direction
$\times$ 3 fog levels $\times$ 3 times of day $\times$ dry/wet road
$\times$ occlusion), each run 3 times. One fog level and one time of day lie
outside the assumptions by design, so 160 of the 288 scenarios break at
least one assumption and their failures must be judged spurious. This gives
the validity check a ground truth: every simulation's verdict
(held/broken) must match its scenario's grammar levels.

\textbf{Pre-computed follow-ups.} The same 288 scenarios, with the same
seeds, are also run for three variants of the system: a path-aware braking
controller, a larger detector (YOLOv5m), and both. A participant who mitigates
at the system layer therefore sees the result of their choice on identical
scenarios. Decisions at other layers are re-judged on the stored traces, or
recorded when they fall outside the four variants.

\textbf{Procedure.} \todo{N} participants with software-engineering, ML or
automotive background each review the same pre-computed round independently,
after a fixed briefing and a guided tour on a training round. The tool's
suggested verdicts are hidden. \toolname{} records every interaction per
participant: \texttt{decisions.json}, an append-only time-stamped log of
verdicts, changes, evidence opened and time per decision, and the signed
$R_1$.

\textbf{Measures.} Per participant: (M1) obstacles accepted, renamed,
rejected and deferred; (M2) requirements and domain assumptions refined, by
label [S]/[D]/[R]/[T] and mitigation layer; (M3) diagnoses reported as
unhelpful; (M4) the before/after result of their revised requirement.
Across participants: where decisions on the same obstacle converge and where
they diverge.

\textbf{Results.} \todo{From the study: M1--M4, the divergence analysis, and
the validity check on the pre-computed space.}

\textbf{Early evidence.} In the authors' own use, the loop shifted a failure
mode rather than removing it. In round~1, emergency braking failed 91 of 120
encounters (76\%). Switching to proportional braking ([S]) reduced failures
to 4 of 128 (3\%), but the car now stopped in 113 of its passes and never
drove on. The per-step traces also exposed defects in our own test scene
(e.g.\ a pedestrian who never crossed), which \toolname{} labels [T] rather
than blaming the system.

\section{Conclusion}
\label{sec:conclusion}
\toolname{} turns REAL's failure-driven loop into a tool. The requirement
keeps the system, the domain and the requirement apart, and the tool checks
the domain assumptions for every simulation. It groups valid failures into
obstacles with evidence, and records the engineer's decisions as a labelled,
signed revision linked to the runs behind it. Next, we will report the study
and extend the system variants beyond braking and perception.
